% New standalone increment. All prior files and GitHub remain read-only.
\documentclass[11pt]{article}
\usepackage[T1]{fontenc}\usepackage[utf8]{inputenc}\usepackage{lmodern}
\usepackage[margin=.88in]{geometry}
\usepackage{amsmath,amssymb,amsthm,mathtools,microtype,array,longtable,booktabs,xurl}
\usepackage[hidelinks]{hyperref}
\setlength{\emergencystretch}{3em}\allowdisplaybreaks[2]
\numberwithin{equation}{section}
\newtheorem{theorem}{Theorem}[section]
\newtheorem{lemma}[theorem]{Lemma}\newtheorem{proposition}[theorem]{Proposition}
\theoremstyle{remark}\newtheorem{remark}[theorem]{Remark}
\newcommand{\E}{\mathbb E}\newcommand{\dd}{\,d}\newcommand{\one}{\mathbf1}
\newcommand{\KK}{\mathcal K}\newcommand{\BB}{\mathcal B}
\newcommand{\Gp}{{G_+}}\newcommand{\eps}{\varepsilon}
\newcommand{\Kfun}{\mathsf K}\newcommand{\Ffun}{\mathsf F}
\newcommand{\ph}{\phi(0)}
\title{Weak-chart reduction and selection at a fixed logistic level}
\author{Standalone analytical increment for independent review}
\date{Version 1 --- 1 November 2026}
\begin{document}\maketitle
\begin{abstract}
We reduce the original weak-receiver velocities on the reached growing
chart to the value-error form of TAILREP. A tangent Gaussian integration
by parts retains both receiver and donor gates without differentiating
a gate. The resulting errors have admissible linear and cubic growth.
Contraction is taken from a comparator with the actual positive clock
donor measure and the positive core output moment. Fixed incoming
weights need no derivative-weight budget. Together with SRR, the
reduction proves selection at a sufficiently small fixed logistic mass:
the strong-reference error and the incoming-law weak $L^2$ error are
$O(q^{1/20})$. The result stops at that mass level.
\end{abstract}

\section{Step 1: conventions, domain and consumer hypotheses}
\label{an03wcsv1:inventory}
Work in the initialized deep regime $\kappa=10$,
$\rho\in[13/20,7/10]$, $\lambda=\rho^2$, $Q=\lambda+q^2$,
$\ell=\log(1/q)$ and $k=10(1-\lambda)$. All uses of RB, WBR,
MP and SRR below are on their proved original-flow intervals.
Keep RB's $t_0=t_h$ and $t_b=\inf\{c=c_b\}$. Keep the fixed
partition $\KK,\Gp,\Sigma,\BB$, with $\mathcal D=\varnothing$,
and the closed selected set
\[
 S=\mathcal C_q=\{b\in[\pi/2,\pi]:-\cos b\ge q^{1/4}\},
 \qquad P_q=\Gp\setminus S,\qquad G=\Gp\cup\Sigma.
\]
For a label, $U=\sqrt m(\cos\theta,\sin\theta)$ and
$W=\sqrt m(\cos\psi,\sin\psi)$. The original radial measure
is $m\,d\lambda_0$. Set
\[
 E_T=\tfrac12\int_T(U_2^2+W_2^2),\quad
 J_T=\tfrac12\int_T(U_1^2+W_1^2),\quad \mu_T=E_T+J_T.
\]
Physical core mass is $M=\mu_\KK$, not the residual guard constant.
The weak coordinates retain SRR's opposite input/output conventions:
\[
 u=(\pi/2-\theta)/q,\qquad v=(\psi-\pi/2)/q.
\]
Fix SRR's model constants $\delta,g=g_W,C_s=980$, and its fixed
admissible $h<1/4$. Let $t_s=\inf\{H=q^{1/8}\}$ and
$t_e=\inf\{H=h\}$, where $H=H_S$ is the actual clock. Define
\begin{equation}
 N'=\lambda N(1-N),\quad N(t_s)=N_0=q^{1/8},\qquad
 t_n=\inf\{N=n\},\quad 0<n<h.
 \label{an03wcsv1:logistic}
\end{equation}
For every fixed $n<h$, MP proves $t_s<t_n<t_e<t_b$ for small $q$.
The level $n$ is chosen only after the weak constants below.

The unchanged reference is CFA/MOMC's actual characteristic
$z_c=(x_c,y_c)$. Define
\begin{equation}
 E_s=|z_c-(X_*,Y_*)(N)|_2,\qquad
 \eta_s(db)=\frac{m(t_s,b)}{N_S(t_s)}d\lambda_0(b),\qquad
 W_2=\big\||u-U_*|+|v-V_*|\big\|_{L^2(\eta_s)},
 \label{an03wcsv1:objects}
\end{equation}
where $N_S=\mu_S$. The probability $\eta_s$ is transported with
fixed labels. Current radial probability is $\eta_t$, and SRR proves
$\omega=d\eta_t/d\eta_s=1+o(1)$, uniformly including the closed
endpoints. Differentiating the norm in \eqref{an03wcsv1:objects}
differentiates only label velocities. Thus no derivative of $\omega$
occurs; its value is used only inside current donor integrals.

\paragraph{Every hypothesis of the invoked selection implication.}
TAILREP \texttt{selection} requires: exact logistic $N$;
$0<N_0<n<1$ and the terminal hit; nonnegative errors; the strong and
weak differential inequalities up to the first exit $E_s<\delta$;
fixed positive $g,C_s,C_w$; a single pointwise constant $d_0\ge0$;
a nonnegative coefficient $\ell_B$ with a common finite integral
bound $L_*$; separately maintained weak return; its specified
small-gain inequality $\Theta\le1/4$ and initial guard
$K_0(E_0+d_0+A_0\sqrt n)\le\delta/2$. The terminal convergence
uses $N_0,A_0,d_0\to0$. Bounded support for each fixed $q$, or
sufficient moments and absolute continuity, is needed for the norm
calculation; WBR's reached chart supplies bounded support for each $q$.

The WEAK \texttt{Wineq} argument additionally requires a common
leading logistic comparator, a nonnegative donor measure, and a
nonnegative-input comparison segment; its scalar derivative uses
$u+Y\ge0$. Its original statement obtains return from
\texttt{forcing/concentration}, assuming $M=1$, fixed normalized
weak weights, $Y\ge0$, bounded $|\rho K_s|$, initial $u\ge0$,
$0\le Y\le Y_{\max}<2\ph$, and the strict margin
\[
 C_u n+U_0^+\sqrt{N_0n}<\ph-Y_{\max}/2,\qquad
 C_u=\ph+\frac{\ph}{\lambda(1-n)}.
\]
We do not apply that leading return theorem to the original flow.
WBR already gives strict original return $u>0$ on $[t_0,t_e]$,
including endpoints. AN06 gives $0<U_*<1$, $0<V_*<11/10$
for $0<N\le1/2$. Below we prove the required positivity of the
actual core moment, pay $M-1$, and use the derivative calculation
in WEAK's proof with the actual positive donor measure. No assumption
of unit original core mass or fixed \emph{current} radial weights
is made.

\section{Step 2: exact weak receivers and polynomial value errors}
\label{an03wcsv1:fieldsection}
Write $s_u=\sin(qu)$, $c_u=\cos(qu)$, and similarly for $v$.
Let
\[
 a=(s_u,c_u),\quad p=(c_u,-s_u),\quad
 b=(-s_v,c_v),\quad b^\perp=(-c_v,-s_v),\quad A=\{a\cdot X>0\}.
\]
The exact scaled velocities, with the mixture factors displayed, are
\begin{equation}
 2u'=\frac1q\sum_{j=1}^2\E_{P_j}[(b\cdot(X-f))(p\cdot X)\one_A],
 \qquad
 2v'=\frac1q\sum_{j=1}^2\E_{P_j}[(b^\perp\cdot(X-f))(a\cdot X)\one_A].
 \label{an03wcsv1:exact}
\end{equation}
WBR and MP prove, throughout $[t_s,t_e]$,
\begin{equation}
 0<u,\quad |u|+|v|\le L=C_L(1+q^{-p_*}),\quad
 p_*=\frac{603}{640},\quad qL=o(1),\quad
 J_S\le C(q^3+q^2H).
 \label{an03wcsv1:chart}
\end{equation}
In particular $c_u\ge1/2$ and the $P_2$ receiver gate has a
uniform Gaussian tail. No strong-chart representation is applied
at a weak receiver.

\subsection{A tangential extension of the boundary integration by parts}
\begin{lemma}[Exact $P_1$ projected moments]
\label{an03wcsv1:IBP}
On $P_1$, set $T=s_u Z_1+c_u Z_2$, $Z_\perp=c_u Z_1-s_u Z_2$
and $r=s_u/q$. These are independent standard normals,
$a\cdot X=q(r+T)$, $p\cdot X=c_u+qZ_\perp$, and
$A=\{T>-r\}$. Thus
\begin{align}
 \E[X_1(p\cdot X)\one_A]
 &=c_u[(1+q^2)\Phi(r)+qs_u\phi(r)],\notag\\
 q^{-1}\E[X_2(p\cdot X)\one_A]
 &=c_u^2\phi(r)-qs_u\Phi(r),\notag\\
 q^{-1}\E[X_1(a\cdot X)\one_A]
 &=\Kfun(r)+qs_u\Phi(r),\notag\\
 q^{-1}\E[X_2(a\cdot X)\one_A]&=qc_u\Phi(r).
 \label{an03wcsv1:targetmoments}
\end{align}
For any fixed donor output $f_{T_0,i}$,
\begin{equation}
 \E[f_{T_0,i}Z_\perp\one_A]
 =q\int_{T_0}W_i(U\cdot p)
       \Pr\{U\cdot X>0,\ A\}\,d\lambda_0.
 \label{an03wcsv1:tangentstein}
\end{equation}
No receiver gate is differentiated.
\end{lemma}
\begin{proof}
Condition on $T$ and integrate by parts only in $Z_\perp$.
The receiver gate is constant in that variable. The ReLU derivative
is its donor indicator almost everywhere, giving
\eqref{an03wcsv1:tangentstein}; Gaussian domination justifies
integration over labels. At $u=0$, $Z_\perp=Z_1$, recovering WBR's
$X_1$ integration by parts. At nonzero $u$, direct integration in
$X_1$ would differentiate the receiver gate, so the tangent rotation
is essential.

For \eqref{an03wcsv1:targetmoments}, use
$\E(T+r)_+=\Kfun(r)$ and
$\E(T+r)_+^2=(1+r^2)\Phi(r)+r\phi(r)$.
The identities
$X_1=s_u(a\cdot X)+c_u(p\cdot X)$ and
$X_2=c_u(a\cdot X)-s_u(p\cdot X)$ then give all four formulas.
In particular $\E(T+r)_+^2-r\Kfun(r)=\Phi(r)$, which
accounts for the last cancellation.
\end{proof}

For later polynomial bookkeeping, uniformly for $u\ge0$, $qu\le1/2$,
\begin{equation}
 |\Phi(s_u/q)-\Phi(u)|+|\phi(s_u/q)-\phi(u)|\le Cq^2,
 \qquad |\Kfun(s_u/q)-\Kfun(u)|\le Cq^2u^3.
 \label{an03wcsv1:gaussiansaturation}
\end{equation}
Indeed $u/2\le s_u/q\le u$ and $|s_u/q-u|\le q^2u^3/6$.
The mean value theorem and boundedness of
$u^j e^{-u^2/8}$ for fixed $j$ prove the first estimate;
$|\Kfun'|\le1$ proves the second. This analytic Gaussian
saturation matters: the crude Lipschitz estimate for $\Phi$,
multiplied by $v$, would produce the inadmissible quartic
$q^2|v|u^3$. It is replaced by the displayed bounded-kernel
estimate, not discarded.

\subsection{The own-cluster coefficient and the weak functional}
Define the actual core moments and actual positive selected measure by
\[
 Y=\int_\KK my,\quad K_c=\int_\KK m\Kfun(\rho x),\qquad
 \mathcal W_S(u)=\int_S\Ffun(u,\widetilde u)\,d\mu_S,\quad
 \mathcal V_S=\int_S\widetilde v\,d\mu_S.
\]
The physical core coordinates here are $\theta=qx,\psi=qy$.
Put $V_O=-\E_{P_2}[(f_O)_1X_2]/(qQ)$ as in CFA/SRR.

\begin{lemma}[Full-gate $P_2$ values]
\label{an03wcsv1:P2}
The $P_2$ contributions to \eqref{an03wcsv1:exact} are
\begin{align}
 2u'_{P_2}&=-\lambda(1-c)u+\mathcal R_{2u},\notag\\
 2v'_{P_2}&=\rho K_c-\lambda V_O-\lambda(1-c)v+\mathcal R_{2v},
 \qquad
 |\mathcal R_{2u}|+|\mathcal R_{2v}|
 \le C_{\rm fixed}q^2(1+|u|+|v|)^3+Cq^{-j}e^{-c_*/q^2}.
 \label{an03wcsv1:P2values}
\end{align}
Here the subscript means the contribution to the total velocity;
the factors two in \eqref{an03wcsv1:exact} remain in force.
\end{lemma}
\begin{proof}
Removing the receiver gate costs at most $Cq^{-j}e^{-c_*/q^2}$:
$\rho c_u\ge\rho/2$, the total radial mass is bounded by RB,
and $|f(X)|\le\mu_{\rm tot}|X|$. Both sides of the failed gate
are covered by Gaussian absolute moments.

Let $B_{ij}=\E_{P_2}[(X_i-f_i)X_j]$. Exactly,
\[
 B_{22}=Q(1-c),\quad
 B_{11}=q^2\left(1-\int W_1U_1\Pr_{P_2}(U\cdot X>0)\right),
 \quad
 B_{21}=-q^2\int W_2U_1\Pr_{P_2}(U\cdot X>0).
\]
The last two formulas are full-gate Stein identities.
Core geometry, $J_G=O(q^2)$ and $\mu_\BB\le q^{11/10}$
give $B_{11}=O(q^2)$ and $B_{21}=O(q^3)$.
We restate and prove the compact core moment used in MOMC/CFA.
For a core donor, put $t_x=\tan(qx)/q$ and $V=\rho+qZ_2$.
Conditioning on $V$ gives exactly
\[
 \E_{P_2}[(f_\KK)_1X_2]
 =q\int_\KK m\cos(qy)\cos(qx)\,
                          \E_V[V\Kfun(t_xV)].
\]
On the fixed donor chart, $t_x=x+O(q^2)$ and both cosines
are $1+O(q^2)$. The second derivative in $V$ of
$V\Kfun(t_xV)$ has a uniform polynomial bound. Taylor expansion,
$\E(V-\rho)=0$, and $\E(V-\rho)^2=q^2$ therefore give
$q\rho K_c+O(q^3)$, using bounded core mass. Hence
\[
 B_{12}=-q\rho K_c+qQ V_O+O(q^3).
\]
This core error is uniform on the proved fixed donor chart.
SRR's balance estimate bounds $V_O$ uniformly.
The two exact contractions of this matrix are
\[
 q^{-1}[-s_vc_uB_{11}+c_vc_uB_{21}+s_vs_uB_{12}-c_vs_uB_{22}],
\]
\[
 q^{-1}[-c_vs_uB_{11}-s_vs_uB_{21}-c_vc_uB_{12}-s_vc_uB_{22}].
\]
Expand the trigonometric coefficients for $q(|u|+|v|)\le1/2$.
The terms from $B_{22}$ give the stated rates, including the
$O(q^2)$ replacement of $Q$ by $\lambda$.
The cubic sine error is retained; all other errors have degree at
most three. This proves \eqref{an03wcsv1:P2values}.
\end{proof}

\begin{lemma}[Finite-noise weak overlap]
\label{an03wcsv1:overlap}
Set $\mathcal W_S^q(u)=q^{-1}\E_{P_1}[(f_S)_2\one_A]$.
On the reached chart,
\begin{equation}
 |\mathcal W_S^q(u)-\mathcal W_S(u)|
 \le C_{\rm fixed}\{q+q^2L+q|u|+q^2|u|^3\}.
 \label{an03wcsv1:overlaperror}
\end{equation}
For $T_0=P_q,\Sigma$, the corresponding unsigned functional
has the receiver-uniform bound
\begin{equation}
 q^{-1}\E_{P_1}[|(f_{T_0})_2|\one_A]
 \le C(E_{T_0}+q^{-1}\sqrt{E_{T_0}J_{T_0}}).
 \label{an03wcsv1:omittedfunctional}
\end{equation}
\end{lemma}
\begin{proof}
The exact selected integrand is
\[
 m\cos(q\widetilde v)\,
 \E\left[\left(\frac{\sin(q\widetilde u)}q
       +\sin(q\widetilde u)Z_1+\cos(q\widetilde u)Z_2\right)_+
 \one_{\{Z_2>-\tan(qu)/q-\tan(qu)Z_1\}}\right].
\]
Replace the positive part by $(Z_2+\widetilde u)_+$.
Its $L^1$ error, before label integration, is at most
$C[q|\widetilde u|+q^2|\widetilde u|^2+q^2|\widetilde u|^3]$.
Replacing $\cos(q\widetilde v)$ by one costs additionally
$Cq^2\widetilde v^2(1+|\widetilde u|)$.
MP's transverse theorem gives
\begin{equation}
 \int_Sm(1+|\widetilde u|+|\widetilde v|
                    +\widetilde u^2+\widetilde v^2)\le C_{\rm fixed},
 \label{an03wcsv1:donormoments}
\end{equation}
since the second moment is at most $CJ_S/q^2=O(q+H)$.
All donor cubic integrals are at most $CL$ times this bound.
Thus these two replacements cost $C(q+q^2L)$.

For the remaining gate replacement condition on $Z_1$.
The integral of $(z+\widetilde u)_+\phi(z)$ between two thresholds
is at most $C(1+|\widetilde u|)$ times their separation.
Here its expected separation is bounded by
$C(q|u|+q^2|u|^3)$. Use \eqref{an03wcsv1:donormoments}.
The resulting exact one-dimensional integral is
$\Ffun(u,\widetilde u)$. This proves \eqref{an03wcsv1:overlaperror}
without a correlation differentiability assumption at coincident
donors. Finally
$|(f_{T_0})_2|\le2\sqrt{E_{T_0}J_{T_0}}|X_1|+2E_{T_0}|X_2|$
proves \eqref{an03wcsv1:omittedfunctional} for every receiver gate.
\end{proof}

\begin{proposition}[Original growing-chart value reduction]
\label{an03wcsv1:reduction}
For every selected receiver on $[t_s,t_e]$,
\begin{align}
 2u'&=\phi(u)-\mathcal W_S(u)-Y\Phi(u)
       -(1-M)v\Phi(u)-\lambda(1-c)u+\mathcal R_u,\notag\\
 2v'&=\rho K_c-\lambda\mathcal V_S-\lambda(1-c)v
                         -(1-M)\Kfun(u)+\mathcal R_v,
 \label{an03wcsv1:leading}\\
 |\mathcal R_u|+|\mathcal R_v|
 &\le C_{\rm fixed}\{q^{1/10}+q(1+|u|+|v|)
                              +q^2(1+|u|+|v|)^3\}.
 \label{an03wcsv1:Rgrowth}
\end{align}
The constant is fixed before $n$; it may depend on the already fixed
passage and level constants. These are value estimates only.
\end{proposition}
\begin{proof}
For the core on $P_1$, remove its donor gate at exponential cost
and set $T_{ij}=\int_\KK W_iU_j$. The proved donor chart gives
\[
 T_{11}=M+O(q^2),\quad T_{21}=qY+O(q^3),\quad
 T_{12}=O(q),\quad T_{22}=O(q^2).
\]
Substitute the four exact moments \eqref{an03wcsv1:targetmoments}
into \eqref{an03wcsv1:exact} for the target minus core.
Equations \eqref{an03wcsv1:gaussiansaturation} and
$|\sin(qv)/q-v|\le Cq^2|v|^3$ give respectively
\[
 \phi(u)-Y\Phi(u)-(1-M)v\Phi(u),\qquad
 -(1-M)\Kfun(u),
\]
with error at most $Cq^2(1+|u|+|v|)^3$.
For explicit bookkeeping: the $X_1p$ moment differs from $\Phi(u)$
by $Cq^2(1+u^2)$; the $X_2p/q$ moment differs from $\phi(u)$
by that same bound; $X_1a/q-\Kfun(u)$ costs $Cq^2(1+u^3)$;
and $X_2a/q$ is $O(q)$. The $T_{12},T_{22}$ factors and
receiver trigonometric factors therefore create no term of degree
larger than three at order $q^2$. Higher products are absorbed using
$q^2L^2=o(1)$. In particular the spurious quartic mentioned after
\eqref{an03wcsv1:gaussiansaturation} is not used.

For $G$ on $P_1$, balance gives
\[
 |(f_G)_1|\le2J_G|X_1|+2\sqrt{E_GJ_G}|X_2|,\qquad
 |(f_G)_2|\le2\sqrt{E_GJ_G}|X_1|+2E_G|X_2|.
\]
Here $J_G=O(q^2)$ and $E_G=O(1)$.
In $2u'$, the $f_{G,1}$ term costs $Cq^2|v|$.
For the second output use \eqref{an03wcsv1:tangentstein}:
\[
 q^{-1}\E[(f_G)_2(p\cdot X)\one_A]
 =c_u\mathcal W_G^q(u)
       +q\int_GW_2(U\cdot p)\Pr\{U\cdot X>0,A\},
\]
whose last term is $O(q^2(1+|u|))$.
The factor $c_vc_u-1$ on the uniformly bounded functional costs
$Cq^2(u^2+v^2)$. The $P_1$ terms in $2v'$ cost
$Cq^2(1+|u|)(1+|v|)$ by the two output bounds and
$|a\cdot X|/q\le C(|u|+|Z_1|+|Z_2|)$.
For $P_q$ and $\Sigma$, \eqref{an03wcsv1:omittedfunctional}
costs respectively $O(q^{1/8})$ and $O(q+q^{\Gamma_\nu/2})$:
use SRR's proved mass/transverse charges, retaining the $q^2$
term in $\mu_{P_q}$. No uncharted mass is divided by $q$ alone.
The $S$ term is Lemma \ref{an03wcsv1:overlap}.

PROFILE's matrix-value estimate bounds the entire $\BB$ field
at any receiver by $C\mu_\BB/q\le Cq^{1/10}$, irrespective of
which side of a donor gate it occupies. This includes both strips.
In the $P_2$ calculation its formal moment is also bounded by
the same amount.
Lemma \ref{an03wcsv1:P2} supplies the own-cluster terms.
Finally SRR \texttt{finiteweak} gives
$|V_S-\mathcal V_S|\le Cq^{37/320}$, while its non-selected
moment charges give $|V_O-V_S|\le Cq^{1/10}$.
All these additive exponents exceed or equal $1/10$.
Also $q^2L=O(q^{677/640})=O(q)$.
This proves \eqref{an03wcsv1:leading}--\eqref{an03wcsv1:Rgrowth}
with every omitted population accounted for.
\end{proof}

\paragraph{Admissibility audit.}
If $d=|u-U_*|+|v-V_*|$, bounded selected coordinates imply
$(1+|u|+|v|)^3\le C[1+(1+L^2)d]$. Thus the growth lines
are split as follows:
\[
\begin{array}{c|c|c|c}
\text{term}&(j,i)&j-p_*(i-1)&\text{coefficient integral}\\ \hline
q(1+|u|+|v|)&(1,1)&1&O(q\ell)\\
q^2(1+|u|+|v|)^2&(2,2)&677/640&O(q^{677/640}\ell)\\
q^2(1+|u|+|v|)^3&(2,3)&37/320&O(q^{37/320}\ell).
\end{array}
\]
Each is admissible. Constants times $q^{1/10}$ are pointwise
additive errors. The exact $-\lambda(1-c)(u,v)$ and physical
mass terms were computed, not bounded as arbitrary large errors.

\section{Step 3: signs, strong moments and donor comparison}
\label{an03wcsv1:donorsection}
RB gives $y\ge9/320$ on every core label and $M\ge1/3$.
Therefore the actual comparison moment satisfies
\begin{equation}
 Y=\int_\KK my\ge3/320>0,\qquad |Y|+|\rho K_c|\le C_{\rm model}.
 \label{an03wcsv1:Ysign}
\end{equation}
If the unexpanded moment $q^{-1}\int_\KK W_2U_1$ is used
instead, it differs from $Y$ by $O(q^2)$ and is at least
$3/640$ for small $q$. Non-core strong-output contributions
have not been silently included in this positive moment:
the $G$ first-output terms are the explicitly charged transverse
terms in Proposition \ref{an03wcsv1:reduction}, and $\BB$ is
charged in absolute value. No sign is asserted for these non-core
fields. They stay in the explicitly budgeted remainder; the
comparator uses the actual core moment itself.

SRR \texttt{corecollapse}, MP's physical mass budget, and the
unchanged reference give
\begin{equation}
 |Y-Y_*|+\rho|K_c-\Kfun(\rho X_*)|
 \le (1+\lambda)M E_s+(1+\lambda)M\bar D_\KK
                +C_{\rm model}|1-M|
 \le3E_s+C_{\rm fixed}q^{1/20}.
 \label{an03wcsv1:strongmoments}
\end{equation}
Here $M\le2$, $\lambda\le49/100$ and
$\bar D_\KK=O(q^{1/20})$; this inequality only compares moments
and does not alter SRR's strong equation.

Let $d_{\rm ph}$ be MP's uniform bound on $|\log(H/N)|$, and set
\begin{equation}
 \epsilon_w(t)=\left[\frac{N_S(t)}{N(t)}
                       \|\omega(t,\cdot)\|_\infty-1\right]_+.
 \label{an03wcsv1:weightcoefficient}
\end{equation}
Then, uniformly on $[t_s,t_e]$,
\begin{equation}
 \epsilon_w\le C[D_q+q^2(1+L^2)+d_{\rm ph}],\quad
 |N_S-N|\le CN(d_{\rm ph}+q^2),\quad
 D_q=C(q^{1/8}+q^2\ell).
 \label{an03wcsv1:weights}
\end{equation}
These are exactly SRR's weight-value conversion and MP's mass
and phase statements. At the bounded comparator input,
\begin{align}
 |\mathcal W_S(U_*)-N\Ffun(U_*,U_*)|
 &\le\int_S|\widetilde u-U_*|\,d\mu_S
                      +C_{\rm model}|N_S-N|,\notag\\
 |\mathcal V_S-NV_*|
 &\le\int_S|\widetilde v-V_*|\,d\mu_S
                      +C_{\rm model}|N_S-N|.
 \label{an03wcsv1:donorcomparison}
\end{align}
The first inequality uses $|\Ffun_r|\le1$ for every donor;
the second is exact linear moment subtraction.
They retain the coefficient one, unlike a coarser combined moment
bound. In particular the weighted sum needed below is at most
\begin{equation}
 \int_S(|\widetilde u-U_*|+\lambda|\widetilde v-V_*|)\,d\mu_S
 \le N(1+\epsilon_w)W_1(\eta_s)
 \le N(1+\epsilon_w)W_2.
 \label{an03wcsv1:coefficientone}
\end{equation}
No derivative-weight term is needed or claimed.

\section{Step 4: the original weak error inequality}
\begin{theorem}[Value-only weak error reduction]
\label{an03wcsv1:weakerror}
One can fix the model-only constant $C_w=20$, independent of $n$,
and a pointwise constant $d_0^{(w)}=C_{\rm wcs}q^{1/20}$,
with $C_{\rm wcs}$ depending only on already fixed constants,
such that up to the first exit from $E_s<\delta$,
\begin{equation}
 D^+W_2\le[-\lambda(1-N)/2+N/2+\ell_B(t)]W_2
                         +C_w(E_s+d_0^{(w)}).
 \label{an03wcsv1:weakineq}
\end{equation}
A valid nonnegative coefficient is
\begin{equation}
 \ell_B=\frac{\lambda}{2}|c-N|+\frac N2\epsilon_w
       +C_{\rm fixed}\{|1-M|+q+q^2(1+L^2)\}.
 \label{an03wcsv1:ellB}
\end{equation}
For small $q$, uniformly over all $n\in(N_0,h)$,
\begin{equation}
 \int_{t_s}^{t_n}\ell_B\le L_*=1.
 \label{an03wcsv1:Lstar}
\end{equation}
In fact the integral tends to zero uniformly.
\end{theorem}
\begin{proof}
Use the comparator scalar input field
\[
 A_t(u)=\tfrac12[\phi(u)-\mathcal W_S(u)-Y\Phi(u)
                                      -\lambda(1-c)u].
\]
Its donors are the actual positive radial measure on $S$.
For every $u\ge0$,
\[
 A_t'(u)=-\tfrac12[(u+Y)\phi(u)
       +\phi(u)\int_S(\widetilde u-u)_+\,d\mu_S+\lambda(1-c)]
       \le-\lambda(1-c)/2.
\]
The segment from the actual $u>0$ to $U_*>0$ lies in this
domain, by WBR and AN06. This is the derivative of the comparator,
not a receiver derivative estimate for a finite-noise remainder.
Use the corresponding output comparator
$[\rho K_c-\lambda\mathcal V_S-\lambda(1-c)v]/2$,
whose scalar derivative is $-\lambda(1-c)/2$.

Subtract the selected branch equations at the exact logistic mass.
The input donor difference and the output linear donor difference
are \eqref{an03wcsv1:donorcomparison}; their sum is controlled by
\eqref{an03wcsv1:coefficientone}.
Strong moment differences cost half of
\eqref{an03wcsv1:strongmoments}.
Changing $c$ to $N$ costs
$\lambda|c-N|d/2$ plus a bounded-comparator
$C_{\rm model}|c-N|$.
The two physical mass terms in \eqref{an03wcsv1:leading} cost
$C_{\rm model}|1-M|(1+d)$ since $\Phi\le1$ and
$\Kfun(u)\le\ph+|u|$.
Finally \eqref{an03wcsv1:Rgrowth} gives an additive
$C_{\rm fixed}q^{1/10}$ and coefficient
$C_{\rm fixed}[q+q^2(1+L^2)]d$.
The bounded mass-normalization errors in
\eqref{an03wcsv1:donorcomparison} and the pointwise phase error
are $O(d_{\rm ph}+q^2)$.
Here $d_{\rm ph}=O(q^{1/8}+q^2\ell)=O(q^{1/8})$ for the fixed
$h$, and $|1-M|=O(q^{11/10})$.
All additive errors other than the core moment have exponent
at least $1/10>1/20$. Thus $C_{\rm wcs}$ can be fixed to put
them in $20d_0^{(w)}$. The coefficient of $E_s$ is at most
$3/2$; choosing $C_w=20$ is model-only.

The resulting common labelwise inequality for
$d=|u-U_*|+|v-V_*|$ has the common donor term
$N(1+\epsilon_w)W_1/2$. Integrate its scalar comparison against
the fixed $\eta_s$, using Minkowski and $W_1\le W_2$.
This proves \eqref{an03wcsv1:weakineq} in the integrated sense
and hence for the upper Dini derivative where needed.
The argument covers zeros of either coordinate difference and
zeros of the norm by the same integrated comparison.
The coefficients are continuous or locally bounded on the finite
chart; the original characteristics and these norms are absolutely
continuous. No current-weight differentiation occurs.

For the coefficient integral, MP gives
\[
 \int_{t_s}^{t_e}|c-N|
 \le C\left(d_{\rm ph}h/b_0+q^{1/4}h/b_0+q^2\ell\right)=o(1),
 \qquad b_0=Q(1-c_b).
 \]
Exact logistic integration and $N(t_e)<2h<1/2$ give
$\int_{t_s}^{t_e}N\le\lambda^{-1}\log(1/(1-2h))$.
MP bounds $\int|1-M|=O(q^{11/10}\ell)$, and RB's duration
is at most $20\ell$. Hence
\begin{align}
 \int_{t_s}^{t_e}\ell_B
 \le C_{\rm fixed}\big[
 &d_{\rm ph}h/b_0+q^{1/4}h/b_0+q^2\ell\notag\\
 &+(D_q+d_{\rm ph}+q^2(1+L^2))
                      \lambda^{-1}\log(1/(1-2h))\notag\\
 &+(q^{11/10}+q+q^{37/320})\ell\big]=o(1).
 \label{an03wcsv1:integratedledger}
\end{align}
The right side uses no terminal $n$, proving the uniform choice
$L_*=1$. This is a proved coefficient budget; no integrated
quantity has been substituted for $d_0^{(w)}$.
\end{proof}

\section{Step 5: a fixed level and the selection theorem}
\label{an03wcsv1:selectionsection}
Let $\lambda_{\min}=169/400$, $g_{\min}=1183/41600$,
$C_s=980$, $C_w=20$, $L_*=1$, and
$B_{\max}=e\,2^{1/(2\lambda_{\min})}$. Fix, for example,
\begin{equation}
 n=\frac12\min\left\{h,\frac14,
              \frac{\lambda_{\min}g_{\min}}{64C_sB_{\max}C_w}\right\}.
 \label{an03wcsv1:nchoice}
\end{equation}
This is a positive fixed number strictly below $h$.
With the actual model rate $g=g_W$, define exactly as TAILREP
\begin{gather}
 B_*=e(1-n)^{-1/(2\lambda)},\qquad
 a_0=\lambda(1-n)/2,\qquad m_0=\min(g/2,a_0/2),\notag\\
 L_0=1+\frac{C_s(1+B_*)}{g-m_0},\qquad K_0=2L_0,\qquad
 \Theta=\frac{C_sB_*C_wn}{(a_0-m_0)(g-m_0)}.
 \label{an03wcsv1:selectionconstants}
\end{gather}
For $n\le1/2$, $B_*\le B_{\max}$,
$a_0-m_0\ge\lambda_{\min}/8$ and
$g-m_0\ge g_{\min}/2$. Thus
\[
 \Theta\le\frac{16C_sB_{\max}C_wn}{\lambda_{\min}g_{\min}}
 \le\frac14.
\]
The half in \eqref{an03wcsv1:nchoice} leaves additional strict
slack. No later value of $q$ or $n$ chooses $C_w$ or $L_*$.

\begin{theorem}[Initialized $L^2$ selection inside the weak passage]
\label{an03wcsv1:main}
Fix constants in the order
\[
 \mathrm{WPT/RB}\ \longrightarrow\ \mathrm{WBR}\
 \longrightarrow\ (a_s=1/8,p=2)\
 \longrightarrow\ (\delta,g_W,C_s)\
 \longrightarrow\ h\ \longrightarrow\ (C_w,L_*)\
 \longrightarrow\ n\ \longrightarrow\ q,
\]
with SRR's unchanged local constants and
\eqref{an03wcsv1:nchoice}--\eqref{an03wcsv1:selectionconstants}.
For the original initialized flow and all sufficiently small $q$,
the strong-reference guard closes on $[t_s,t_n]$, and
\begin{equation}
 \begin{gathered}
 E_s(t)\le K_0[E_0e^{-m_0(t-t_s)}+d_0+A_0\sqrt{N(t)}],\\
 d_0=\max\{C_{\rm srr},C_{\rm wcs}\}q^{1/20},\qquad
 E_0=O(q^{1/20}),\qquad A_0=O(q^{1/16}).
 \end{gathered}
 \label{an03wcsv1:Econclusion}
\end{equation}
Writing $T_n=t_n-t_s$,
\begin{align}
 W_2(t_n)\le B_*\left[\frac{A_0}{\sqrt n}
 +C_w(K_0+1)\left\{
 \frac{E_0e^{-m_0T_n}}{a_0-m_0}+\frac{d_0}{a_0}
                           +\frac{A_0\sqrt n}{2a_0}\right\}\right].
 \label{an03wcsv1:Wconclusion}
\end{align}
In particular
\begin{equation}
 E_s(t_n)+W_2(t_n)=O_n(q^{1/20}).
 \label{an03wcsv1:rates}
\end{equation}
This is selection in the fixed-incoming-law $L^2$ norm, not a
uniform labelwise assertion. No conclusion past $t_n$ is claimed.
\end{theorem}
\begin{proof}
The exact logistic solution gives $0<N_0=q^{1/8}<n<1$ and the
terminal hit, with $t_n<t_e<t_b$ by MP.
WBR maintains original return throughout this interval; AN06
supplies the selected positive branch in $N<1/2$.
The actual positive donor measure and positive $Y$ are verified
above. SRR supplies the unchanged strong inequality up to the
first exit $E_s<\delta$ with pointwise $d_0^{(s)}$.
Theorem \ref{an03wcsv1:weakerror} supplies the other inequality
on precisely that scope, with $\ell_B\ge0$ and
$\int\ell_B\le1$. Enlarging both additive constants to the same
$d_0$ preserves both inequalities.
SRR proves $E_0=O(q^{1/20})$ and MP proves
$A_0=O(q^{1/16})$ for exactly \eqref{an03wcsv1:objects}.
Consequently
$K_0(E_0+d_0+A_0\sqrt n)\le\delta/2$ for small $q$.
The initial state is strictly inside the guard, every constant is
fixed, and both errors are nonnegative and absolutely continuous.
All hypotheses of TAILREP \texttt{selection} are now discharged.
Its conclusions are exactly
\eqref{an03wcsv1:Econclusion}--\eqref{an03wcsv1:Wconclusion}.

For explicit rates, logistic integration yields
\begin{equation}
 T_n=\frac1\lambda\log\frac{n(1-q^{1/8})}{q^{1/8}(1-n)}
 =\frac{\ell}{8\lambda}+\frac1\lambda\log\frac n{1-n}+o(1).
 \label{an03wcsv1:duration}
\end{equation}
Hence $E_0e^{-m_0T_n}=O_n(q^{1/20+m_0/(8\lambda)})$.
The other terms are $O_n(q^{1/20})$ and $O_n(q^{1/16})$,
which proves \eqref{an03wcsv1:rates}. In particular $T_n\to\infty$.
The local strong guard is improved to $E_s\le\delta/2$ by the
invoked theorem, not presumed during the original-field reduction.
\end{proof}

\section{Step 6: the S5 handoff, with accuracy stated}
\label{an03wcsv1:handoff}
All constants in this table are fixed after the displayed constant
order. Rates with subscript $n$ permit dependence on the chosen
fixed positive level.
\begin{longtable}{@{}p{.27\textwidth}p{.67\textwidth}@{}}
\toprule
Quantity at $t_n$ & Proved state\\ \midrule\endhead
Actual strong reference &
$|z_c-Z_*(n)|_2=O_n(q^{1/20})$, and $E_s\le\delta/2$
through the selection interval.\\
Core first moment and mass &
$\int_\KK m(|x-X_*(n)|+|y-Y_*(n)|)=O_n(q^{1/20})$;
$\mu_\KK=1+O(q^{11/10})$; inherited radius
$r_\KK=O(q^{1/20})$.\\
Weak distance &
$W_2(\eta_s)=O_n(q^{1/20})$ to $(U_*,V_*)(n)$.
Under current weights,
$W_2(\eta_{t_n})\le\|\omega\|_\infty^{1/2}W_2(\eta_s)
=(1+o(1))W_2(\eta_s)$.\\
Selected mass and response &
$N_S/n=1+O_n(q^{1/8})$,
$c(t_n)=n+O_n(q^{1/8})$.
The actual clock is $H(t_n)=n\exp(O(d_{\rm ph}))$.\\
Non-selected mass &
$\mu_{P_q}+\mu_\Sigma+\mu_\BB
\le C(q^{1/4}H+q^2)=O_n(q^{1/4})$.
More specifically
$\mu_\Sigma\le C(q^5+q^{\Gamma_\nu}H+q^2)$
and $\mu_\BB\le q^{11/10}$.\\
Non-selected fields &
The omitted weak nonlinear functionals and linear moments cost
$O(q^{1/8})$ for $P_q$,
$O(q+q^{\Gamma_\nu/2})$ for $\Sigma$,
and $O(q^{1/10})$ for $\BB$.
At actual receivers the remaining transverse terms keep their
explicit admissible polynomial factors.\\
Outside budgets &
RB's proved pointwise and integrated ancestry budgets remain valid
on this subinterval. In particular $\sup\mu_\BB/q\le q^{1/10}$
and $\int_{t_0}^{t_n}\mu_\BB/q\le20q^{1/10}\ell=o(1)$.\\
Phase/time &
With the exact entry clock $H_0=H(t_0)$,
$t_n-t_0=Q^{-1}\log\{n/[(1-n)H_0]\}+o(1)$.
Since $H_0\asymp A_q^{-\lambda}q^k$, its expression using
only $A_q$ has an $O(1)$ seed uncertainty.\\
\bottomrule
\end{longtable}
The first-moment statement follows from SRR's core radius and
\eqref{an03wcsv1:rates}; fixed mass multiplies the reference
error. The mass and response statements use MP, not a derivative
of a small phase error.

For completeness, the time accuracy in the table follows from RB's
exact identity
\[
 t_n-t_0=\frac{\log(H(t_n)/H_0)-\delta_H(t_n)}Q
                         +\int_{t_0}^{t_n}c,\qquad
 |\delta_H|\le D_q=o(1).
\]
Before $t_s$, SRR \texttt{preforcing} gives
$c\le C(H+q^2+\mu_\BB)$. The proved backward clock bounds
$\int_{t_0}^{t_s}H\le Cq^{1/8}$, so that part of $\int c$
is $o(1)$, retaining the outside term.
After $t_s$, MP gives $\int|c-N|=o(1)$ and
$\int_{t_s}^{t_n}N=\lambda^{-1}
\log[(1-N_0)/(1-n)]$. Combine these with
$H(t_n)/n=e^{O(d_{\rm ph})}$.
Replacing $\lambda$ by $Q$ only in the fixed logarithmic
factor costs $O(q^2)$, proving the asserted $o(1)$ phase.
Equivalently, with only the entry asymptotics,
\[
 t_n-t_0=
 \frac{k\ell+\lambda\log A_q+\log(n/(1-n))+O(1)}Q.
\]
This does not use RB's terminal time as an $o(1)$ clock.

\paragraph{Limitations.}
The result is $p=2$ selection at a fixed small level and a
population-to-orbit value comparison. It supplies neither
$p=\infty$ selection nor pairwise finite-noise contraction.
The growing chart remains a valid support bound; the $L^2$
conclusion does not replace it by uniform bounded support.
Compact-window transfer, any continuation beyond $t_n$, and all
later retention belong to S5. No prior increment is modified.

\begin{thebibliography}{12}\small\raggedright
\bibitem{RB} RB, \path{AN03_response_barrier_v1.tex}:
\texttt{an03rbv1:completed, conclusions, massconclusions, primaryclock, clock}.
\bibitem{WBR} WBR, \path{AN03_weak_boundary_return_v1.tex}:
\texttt{an03wbrv1:exact, boundaryidentity, main, clockalignment, latechart, mass}.
\bibitem{MP} MP, \path{AN03_matching_package_v1.tex}:
\texttt{an03mpv1:accounting, response, phase, phaseintegral, masses,
transverse, smallness, massbudgets, linearbudget}.
\bibitem{SRR} SRR, \path{AN03_strong_reference_reduction_v1.tex}:
\texttt{an03srrv1:main, W, weights, momentcharge, finiteweak, weakmoments,
corecollapse, delta, hbound, preforcing}.
\bibitem{CFA} CFA, \path{AN03_c_form_assembly_v1.tex}:
compact core Gaussian moments and the exact definition of $V_O$.
Its strong receiver chart is not used for weak receivers.
\bibitem{MOMC} MOMC, \path{AN02_moving_overlap_M_clock_passage_v1.tex}:
the compact core Gaussian moments in its strong reduction proof.
\bibitem{P} P, \path{THEOREM_A_ANALYTICAL_PROGRESS.tex}:
\texttt{pa:F, pa:pop-u, pa:pop-v, pa:5d, pa:branch}.
\bibitem{AN06} AN06, \path{AN06_learned_orbit_witnesses_v2.tex}:
\texttt{inc:AN06:v2:regiontheorem}.
\bibitem{WEAK} WEAK, \path{AN03_weak_transit_forcing_and_orbit_selection_v1.tex}:
\texttt{inc:AN03:weak:v1:forcing, concentration, Wineq}.
\bibitem{TAILREP} TAILREP,
\path{AN03_mass_charged_tail_and_selection_defects_v1.tex}:
\texttt{inc:AN03:tailrepair:v1:weakerror, selection, Econclusion, Wconclusion}.
\bibitem{PROFILE} PROFILE,
\path{AN02_strong_entry_profile_and_tail_budget_v1.tex}:
\texttt{inc:AN02:profile:v1:cost}.
\end{thebibliography}
\end{document}
